\section{Cost accounting}\label{sec:costs}

This section answers the sample-requirement part of Q4 and gives the cost side of the comparison of \cref{sec:compare}.

\subsection{Accounting rules}

The rules were fixed in the specification before any run.
\begin{enumerate}[nosep,label=(\roman*)]
\item \emph{Training transitions} are the environment transitions that a learner receives: the on-policy rollouts of $f$-PG, and the task-blind replay buffer of FB. The four demonstrations ($\vn{s14.demo_transitions}{400}$ transitions, states only) are common to all methods and counted separately.
\item \emph{Evaluation transitions} are fresh rollouts generated by the evaluator. They are never given to a learner and never counted as training data; they are reported separately.
\item FB pretraining is task-blind and serves the three goals of a layout and seed. The \emph{one-task} cost charges the whole pretraining to a single task; the \emph{amortized} cost divides it by the three goals that used it. The adaptation cost of \ERFB{} and \DMFB{} is added in both cases; it uses no new transition.
\item Wall times are measured on one laptop CPU with \vn{s10.threads}{4} PyTorch threads (NumPy and BLAS thread counts were neither set nor recorded), with jobs run one after another. They include the diagnostics that the training loops compute in passing (for example the logged objectives and entropies); this overhead was not timed separately and is not removed. Two conventions are used, never mixed within a comparison. \emph{Per-task phase times} (\cref{tab:costs,tab:cost-components} and the comparisons below them) sum the phase timers in each job's record; for FB pretraining they add the recorded generation time of its replay buffer, which the data-preparation step performed once, before and outside the timed jobs. \emph{Whole-job wall times} (the totals at the end of this section) run from the start to the end of each job: they include job set-up, checkpoints and output writing, but not the replay generation. (This description of the two conventions and of the diagnostics was added after the runs; it changes no measurement.)
\end{enumerate}

\begin{table}[ht]
\centering\small
\caption{Training and evaluation cost per task (core suite; means over the \vn{n.runs.finite}{18} finite and \vn{n.runs.cont}{9} continuous runs of each method). Training transitions are those given to the learner; evaluation transitions are generated by the evaluator only. Wall times are per-task phase times in seconds on one laptop CPU (rule (iv)). \ERFB{} and \DMFB{} share one evaluation job per adaptation (\vn{s14.eval.ckpt.er}{4} \ERFB{} and \vn{s14.eval.ckpt.dm}{10} \DMFB{} checkpoints): its cost is shown once for the pair and counted once in the totals.}\label{tab:costs}
\resizebox{\textwidth}{!}{% GENERATED by report/verify/make_tables.py from the saved run files (runs/jobs/*).
% Cost accounting, core suite (CPU, 4 PyTorch threads): mean over the runs of each environment class (finite: 18 per method; continuous: 9).  Wall times in seconds: per-task phase timers (result.json), FB pretraining including the recorded replay generation.  ER_FB and DM_FB share one evaluation job per adaptation: its cost is printed once for the pair.
% Do not edit by hand: rerun  python3 report/verify/verify_results.py

\begin{tabular}{@{}llrrrrrr@{}}
\toprule
 & & \multicolumn{2}{c}{one task} & \multicolumn{2}{c}{amortized over 3 goals} & \multicolumn{2}{c}{held-out evaluation}\\
\cmidrule(lr){3-4}\cmidrule(lr){5-6}\cmidrule(l){7-8}
maze & method & transitions & wall (s) & transitions & wall (s) & transitions & wall (s)\\
\midrule
finite & \fPG{} (JS) & 50{,}000 & 0.38 & 50{,}000 & 0.38 & 120{,}000 & 0.14\\
 & \DMFB & 50{,}000 & 51.78 & 16{,}667 & 17.62 & \multirow{2}{*}{280{,}000} & \multirow{2}{*}{1.28}\\
 & \ERFB & 50{,}000 & 51.25 & 16{,}667 & 17.08 &  & \\
\addlinespace
continuous & \fPG{} (JS) & 100{,}000 & 16.03 & 100{,}000 & 16.03 & 550{,}000 & 2.64\\
 & \DMFB & 100{,}000 & 100.32 & 33{,}333 & 33.80 & \multirow{2}{*}{700{,}000} & \multirow{2}{*}{3.79}\\
 & \ERFB & 100{,}000 & 99.78 & 33{,}333 & 33.26 &  & \\
\bottomrule
\end{tabular}
}
\end{table}

\begin{table}[ht]
\centering\small
\caption{Cost components per phase (core suite; per-task phase times; mean $\pm$ sample SD over runs where it varies, and over the distinct models for FB pretraining, each model counted once).}\label{tab:cost-components}
% GENERATED by report/verify/make_tables.py from the saved run files (runs/jobs/*).
% Cost components per phase (core suite): per-task phase timers (result.json); mean and sample SD over runs, and for FB pretraining over the distinct models (6 finite, 3 continuous), each counted once.
% Do not edit by hand: rerun  python3 report/verify/verify_results.py

\begin{tabular}{@{}lrr@{}}
\toprule
phase & finite & continuous\\
\midrule
expert demonstrations (transitions, every method) & 400 & 400\\
FB task-blind replay (transitions, once per layout and seed) & 50{,}000 & 100{,}000\\
FB replay resets & 1{,}000 & 2{,}000\\
FB updates (batch 128 + target batch 128) & 10{,}000 & 20{,}000\\
FB pretraining wall time per model (s) & $51.2\pm1.0$ & $99.8\pm0.5$\\
\ERFB{} latent at the final snapshot (s) & $0.0005$ & $0.0004$\\
\DMFB{} 300 latent steps at the final snapshot (s) & $0.54\pm0.02$ & $0.54\pm0.01$\\
new transitions during latent adaptation & 0 & 0\\
\fPG{} learner transitions per task & 50{,}000 & 100{,}000\\
\fPG{} policy updates per task & 25 & 50\\
\fPG{} wall time per task (s) & $0.38\pm0.01$ & $16.0\pm0.1$\\
\bottomrule
\end{tabular}

\end{table}

\subsection{What the runs cost}

\Cref{tab:costs,tab:cost-components} give the costs.
By design the two families had the same transition budget for one task: $\vn{budget.same.f}{50{,}000}$ transitions in a finite maze and $\vn{budget.same.c}{100{,}000}$ in the continuous maze, collected on-policy by $f$-PG and task-blind by FB.
Amortized over its three goals, FB used a third of that per task.
The compute differs much more than the data.
In the finite mazes, tabular $f$-PG trained in $\vn{s14.fpg.wall.f}{0.38}$\,s per task, while FB pretraining took $\vn{s14.fb.wall.f}{51.2}$\,s; in the continuous maze, $f$-PG with an MLP policy and a KDE signal took $\vn{s14.fpg.wall.c}{16.0}$\,s and FB pretraining $\vn{s14.fb.wall.c}{99.8}$\,s.
Once a model is trained, imitation is almost free: \ERFB{} needs one weighted average of $B$ over the demonstrations (below \vn{s14.er.ms}{1}\,ms), and the 300 \DMFB{} steps take $\vn{s14.dm.wall}{0.54}$\,s.

\paragraph{Samples to reach the thresholds.}
In the finite mazes $f$-PG first met the operational thresholds at the $\vn{abs.fpg.first.lo}{20{,}000}$- or $\vn{abs.fpg.first.hi}{30{,}000}$-transition checkpoint (\cref{tab:convergence}), i.e.\ within $\vn{s14.pct.lo}{40}$ to $\vn{s14.pct.hi}{60}\%$ of its budget; with checkpoints every $\vn{fpg.interval.f}{10{,}000}$ transitions, these are upper bounds at that resolution.
FB imitation did not reach that quality with its whole budget, except for \vn{abs.er.last}{2} \ERFB{} runs, so the number of task-blind transitions that FB would need to match $f$-PG is not measured: $\vn{replay.f}{50{,}000}$ replay transitions with up to $\vn{fb.updates.f}{10{,}000}$ pretraining updates were not enough in \vn{s14.notlast}{34} of the \vn{s11.curves}{36} finite runs. Across the snapshots only the number of updates varied; the replay size was fixed.
A cost comparison at equal quality is therefore not available, and the amortized figures should not be read as an advantage of FB.
Amortized over its three goals, FB uses fewer task-blind transitions per task (a third of the $f$-PG budget) but more wall time on this CPU: in per-task phase times (\cref{tab:costs}: pretraining with its replay generation divided by three plus the final-snapshot adaptation, against $f$-PG training, evaluation excluded), about $\vn{s14.ratio.f.lo}{45}$ to $\vn{s14.ratio.f.hi}{46}$ times the wall time of $f$-PG in the finite mazes and about $\vn{s14.ratio.c}{2.1}$ times in the continuous maze. And it buys a worse policy.
In the continuous maze neither family reached the thresholds within its budget, so no samples-to-threshold figure exists for either.

\paragraph{Evaluation is the largest consumer of transitions.}
The core suite gave $\vn{s14.fpg.tr}{1{,}800{,}000}$ transitions to the \vn{s10.core.fpg}{27} $f$-PG runs and $\vn{s14.fb.tr}{600{,}000}$ to the \vn{s10.core.fb}{9} FB models, while its evaluations generated $\vn{s14.eval.tr}{19{,}260{,}000}$ transitions, and the ablations another $\vn{s14.abl.eval.tr}{1{,}560{,}000}$.
Evaluating every checkpoint with hundreds of fresh rollouts is a property of this study, not of the methods; in the finite mazes the oracle makes most of it unnecessary, and the rollouts serve as a cross-check.

\paragraph{Total compute.}
In whole-job wall times (rule (iv)), the \vn{s10.core.jobs}{144} core jobs ran for $\vn{s14.wall.total}{932.7}$\,s in total: $\vn{s14.wall.fb}{618.7}$\,s of FB pretraining, $\vn{s14.wall.fpg}{160.2}$\,s of $f$-PG training, $\vn{s14.wall.eval}{90.2}$\,s of evaluation, $\vn{s14.wall.adapt}{61.2}$\,s of latent adaptation and $\vn{s14.wall.evalexp}{2.3}$\,s of expert evaluation. The adaptation jobs compute the \ERFB{} latent and the \DMFB{} steps at all \vn{fb.n_snapshots}{four} FB snapshots for the learning curves, about $\vn{s14.adapt_job}{2.3}$\,s per job; the per-task cost of \cref{tab:costs} counts the final snapshot only. The \vn{s10.abl.new}{21} new ablation jobs took $\vn{s14.abl.wall}{174.4}$\,s. The peak process memory was $\vn{s14.peak}{823}$\,MiB.

\subsection{Where the comparison is not symmetric}

\begin{itemize}[nosep]
\item \textbf{Resets.} FB's replay starts its $\vn{replay.segment}{50}$-transition segments from states drawn uniformly over the free space, a simulator capability that $f$-PG does not use: every $f$-PG rollout starts at $(1,1)$. This favours FB's coverage of the state space.
\item \textbf{Network sizes.} The FB networks are larger than the $f$-PG policies, and finite $f$-PG is tabular. The wall-time ratio reflects these choices and the CPU; the pilot measured the laptop CPU faster than its GPU for these networks, and a different device would change the ratio.
\item \textbf{Tuning.} Both $f$-PG learning rates and the FB temperature were tuned on one pilot seed and goal; the FB network sizes, update counts and \DMFB{} settings were not tuned (\cref{sec:m-pilot}).
\item \textbf{Number of tasks.} With three goals per environment the amortization factor is three. FB's per-task cost falls further with more tasks; $f$-PG's does not. This is the regime where FB imitation is designed to pay off, and the one this study could not reach at equal quality.
\end{itemize}
